% ============================================================================
%  Cuerpo de la presentacion: diapositivas + guion (\note) por slide.
%  Se incluye desde presentacion.tex (limpia) y presentacion-notas.tex (con guion).
% ============================================================================

% ---------------------------------------------------------------------------
%  PORTADA
% ---------------------------------------------------------------------------
{
\setbeamercolor{background canvas}{bg=azultec}
\begin{frame}[plain]
\begin{tikzpicture}[remember picture,overlay]
  % banda diagonal guinda inferior
  \fill[guinda] ([yshift=3.4cm]current page.south west)
    -- (current page.south west) -- (current page.south east)
    -- ([yshift=1.6cm]current page.south east) -- cycle;
  % filo ambar
  \fill[ambar]  ([yshift=3.4cm]current page.south west)
    -- ([yshift=3.15cm]current page.south west)
    -- ([yshift=1.35cm]current page.south east)
    -- ([yshift=1.6cm]current page.south east) -- cycle;
  % eyebrow
  \node[anchor=north west, text=ambar, font=\demi\large]
    at ([xshift=1.15cm,yshift=-1.0cm]current page.north west)
    {INSTITUTO TECNOL\'OGICO DE CIUDAD JU\'AREZ \, $\cdot$ \, TecNM};
  % titulo (una sola linea, sin guion feo)
  \node[anchor=north west, text=white, font=\heavy\fontsize{42}{44}\selectfont]
    at ([xshift=1.05cm,yshift=-1.85cm]current page.north west)
    {Metaverso Escolar};
  % subtitulo
  \node[anchor=north west, text=white, font=\medium\fontsize{17}{21}\selectfont, text width=13cm]
    at ([xshift=1.15cm,yshift=-3.5cm]current page.north west)
    {Pr\'acticas en videojuego conectadas a Moodle,\\ con calificaci\'on autom\'atica};
  % autor / fecha sobre la banda
  \node[anchor=west, text=white, font=\demi\normalsize]
    at ([xshift=1.18cm,yshift=2.15cm]current page.south west)
    {Daniel Octavio Ram\'irez Neri};
  \node[anchor=west, text=white!85, font=\normalsize]
    at ([xshift=1.18cm,yshift=1.4cm]current page.south west)
    {Residencia Profesional \, $\cdot$ \, Junio 2026};
\end{tikzpicture}
\end{frame}
}
\note{
\textbf{Apertura (30 s).} Saluda y presenta: ``Esto es el Metaverso Escolar, un
proyecto que conecta un videojuego de pr\'acticas con Moodle. El alumno juega, y su
calificaci\'on llega sola a la plataforma, sin que el maestro capture nada.''
Aclara que el autor (Daniel) no pudo asistir y que esta presentaci\'on los lleva
paso a paso. Promete una demo al final.
}

% ---------------------------------------------------------------------------
%  EL PROBLEMA
% ---------------------------------------------------------------------------
\begin{frame}{El problema que resolvemos}
\vspace{0.2cm}
\begin{columns}[T]
\column{0.5\textwidth}
\begin{tarjetabox}{rojoitcj}
{\demi\color{rojoitcj}Hoy}\\[2pt]
Las pr\'acticas en simuladores o juegos viven \hl{aparte} de la plataforma escolar.
El maestro \hl{revisa a mano} y \hl{captura calificaciones} una por una.
\end{tarjetabox}
\begin{tarjetabox}{gristenue}
{\demi\color{gristenue}Consecuencia}\\[2pt]
Tiempo perdido, errores de captura y poca evidencia de lo que realmente hizo cada alumno.
\end{tarjetabox}

\column{0.5\textwidth}
\begin{tarjetabox}{verde}
{\demi\color{verde}Con el Metaverso}\\[2pt]
El juego se abre \hlb{desde Moodle} como una actividad m\'as. Al terminar,
la calificaci\'on \hlb{regresa sola} a la plataforma.
\end{tarjetabox}
\begin{tarjetabox}{azulviv}
{\demi\color{azulviv}Resultado}\\[2pt]
Cero captura manual, evidencia automatica del desempe\~no y una experiencia
mucho m\'as atractiva para el alumno.
\end{tarjetabox}
\end{columns}
\end{frame}
\note{
\textbf{El porqu\'e (1 min).} Plantea el dolor: hoy las pr\'acticas digitales no
``hablan'' con Moodle. El maestro tiene que revisar y capturar a mano. Eso cuesta
tiempo y se presta a errores. Recalca la columna verde: con esto, el juego vive
DENTRO de Moodle y la nota se sincroniza sola. Pregunta ret\'orica al p\'ublico:
``\textit{?`Cu\'antas horas se van en capturar calificaciones?}''
}

% ---------------------------------------------------------------------------
%  LA SOLUCION EN UNA FRASE
% ---------------------------------------------------------------------------
\begin{frame}{La idea, en una frase}
\vspace{0.9cm}
\begin{center}
{\heavy\fontsize{26}{32}\selectfont\color{azultec}
El alumno entra a Moodle, abre la actividad,\\[6pt]
\hl{juega la pr\'actica}\, y su \hlb{calificaci\'on aparece sola}\\[6pt]
en el curso.}
\end{center}
\vspace{0.8cm}
\begin{center}
{\medium\color{gristenue}\large Sin instalar nada raro. Sin capturar a mano. Usando el
est\'andar que Moodle ya entiende.}
\end{center}
\end{frame}
\note{
\textbf{Ancla mental (30 s).} Esta es la frase que quieres que recuerden. L\'eela
despacio. El mensaje clave: ``usa un est\'andar que Moodle YA entiende'' (se llama
LTI, lo vemos en un momento) — no es un parche ni un hack, es la forma oficial de
integrar herramientas con Moodle.
}

% ---------------------------------------------------------------------------
%  QUE GANA CADA QUIEN
% ---------------------------------------------------------------------------
\begin{frame}{Qu\'e gana cada quien}
\vspace{0.25cm}
\begin{columns}[T]
\column{0.33\textwidth}
\begin{tarjetabox}{azulviv}
\begin{center}\bolanum{azulviv}{A}\end{center}
\vspace{2pt}{\demi\color{azulviv}\large Alumno}\\[3pt]
Una pr\'actica divertida tipo videojuego. Entra con un clic desde su curso, sin
crear otra cuenta.
\end{tarjetabox}

\column{0.33\textwidth}
\begin{tarjetabox}{guinda}
\begin{center}\bolanum{guinda}{M}\end{center}
\vspace{2pt}{\demi\color{guinda}\large Maestro}\\[3pt]
Califica solo. Ve resultados y un panel propio, sin tocar c\'odigo ni capturar
notas a mano.
\end{tarjetabox}

\column{0.33\textwidth}
\begin{tarjetabox}{verde}
\begin{center}\bolanum{verde}{I}\end{center}
\vspace{2pt}{\demi\color{verde}\large Instituci\'on}\\[3pt]
Tecnolog\'ia est\'andar y abierta. Cada escuela puede tener su propio servidor,
con sus datos.
\end{tarjetabox}
\end{columns}
\end{frame}
\note{
\textbf{Beneficios por p\'ublico (1 min).} Adapta seg\'un qui\'en est\'e en la sala:
si hay directivos, enf\'atiza ``Instituci\'on'' (datos propios, est\'andar, sin
amarre a un proveedor). Si hay docentes, enf\'atiza ``Maestro'' (no capturan nada).
Para todos: el alumno entra con UN clic, sin cuenta nueva.
}

% ---------------------------------------------------------------------------
%  DIVISOR 01
% ---------------------------------------------------------------------------
\divisor{01}{C\'omo se vive}

% ---------------------------------------------------------------------------
%  EL RECORRIDO DEL ALUMNO (linea de tiempo)
% ---------------------------------------------------------------------------
\begin{frame}{El recorrido del alumno, paso a paso}
\vspace{0.7cm}
\begin{center}
\begin{tikzpicture}[every node/.style={font=\small}]
  % linea base
  \draw[line width=2pt, gristenue!50] (0,0) -- (12.4,0);
  \foreach \x/\c/\n in {0/azulviv/1, 3.1/guinda/2, 6.2/azulviv/3, 9.3/guinda/4, 12.4/verde/5}{
    \fill[\c] (\x,0) circle (8pt);
    \node[text=white, font=\cond\small] at (\x,0) {\n};
  }
  \node[text width=2.6cm, align=center, anchor=north] at (0,-0.45)   {\demi Entra a Moodle\\[1pt]\scriptsize con su cuenta de siempre};
  \node[text width=2.6cm, align=center, anchor=north] at (3.1,-0.45) {\demi Abre la actividad\\[1pt]\scriptsize ``Pr\'actica en metaverso''};
  \node[text width=2.6cm, align=center, anchor=north] at (6.2,-0.45) {\demi Se lanza el juego\\[1pt]\scriptsize con un enlace seguro};
  \node[text width=2.6cm, align=center, anchor=north] at (9.3,-0.45) {\demi Realiza la pr\'actica\\[1pt]\scriptsize aciertos y avance};
  \node[text width=2.8cm, align=center, anchor=north] at (12.4,-0.45){\demi\color{verde}Su nota aparece\\[1pt]\scriptsize sola en el curso};
\end{tikzpicture}
\end{center}
\vspace{0.5cm}
\begin{center}
\begin{calloutbox}{azulviv}
\demi El alumno nunca ve nada t\'ecnico: todo pasa entre Moodle, el servidor y el
juego \textbf{de forma autom\'atica}.
\end{calloutbox}
\end{center}
\end{frame}
\note{
\textbf{El flujo completo (1.5 min).} Recorre los 5 c\'irculos con el dedo.
(1) El alumno YA tiene su Moodle. (2) En su curso aparece la actividad como
cualquier tarea. (3) Al darle clic, el juego se abre solo con un enlace seguro de
un solo uso. (4) Juega la pr\'actica. (5) Al terminar, la calificaci\'on viaja de
regreso y aparece en Moodle. Remata con el callout: el alumno no ve nada t\'ecnico.
}

% ---------------------------------------------------------------------------
%  LTI: se ve como una actividad mas
% ---------------------------------------------------------------------------
\begin{frame}{?`Por qu\'e parece parte de Moodle?}
\vspace{0.25cm}
\begin{columns}[T]
\column{0.52\textwidth}
Usa \hl{LTI}, el est\'andar oficial para conectar herramientas con plataformas
educativas. Por eso:
\vspace{4pt}
\begin{itemize}\setlength\itemsep{6pt}
  \item Se agrega como ``\demi herramienta externa'', igual que un Kahoot.
  \item El alumno \hl{no crea otra cuenta}: usa su sesi\'on de Moodle.
  \item Es abierto: \hlb{no amarra} a la escuela con un proveedor.
\end{itemize}
\column{0.48\textwidth}
\begin{tarjetabox}{guinda}
{\demi\color{guinda}En cristiano}\\[2pt]
LTI es el ``cable universal'' entre Moodle y herramientas de fuera. Nuestro juego
se enchufa sin trucos.
\end{tarjetabox}
\begin{tarjetabox}{verde}
{\demi\color{verde}Ya probado}\\[2pt]
Funciona contra un Moodle \textbf{real}, no solo en teor\'ia.
\end{tarjetabox}
\end{columns}
\end{frame}
\note{
\textbf{LTI sin tecnicismos (1 min).} La palabra clave es LTI. No te metas en
siglas: di que es el ``cable universal'' (o ``el USB'') entre Moodle y herramientas
externas. Es el MISMO mecanismo que usan Kahoot, Turnitin, etc. Eso da confianza:
no inventamos nada raro, seguimos el est\'andar que la industria ya usa. Subraya:
ya se prob\'o contra un Moodle real.
}

% ---------------------------------------------------------------------------
%  LA CALIFICACION REGRESA SOLA
% ---------------------------------------------------------------------------
\begin{frame}{La calificaci\'on regresa sola}
\vspace{0.5cm}
\begin{center}
\begin{tikzpicture}[node distance=0.8cm, every node/.style={font=\demi}]
  \node[draw=azulviv, line width=1.5pt, rounded corners=8pt, fill=white,
        text width=2.7cm, align=center, minimum height=1.5cm] (juego)
        {Juego\\[2pt]\scriptsize El alumno termina la pr\'actica};
  \node[draw=azultec, line width=1.5pt, rounded corners=8pt, fill=white,
        text width=2.9cm, align=center, minimum height=1.5cm, right=1.6cm of juego] (server)
        {Servidor\\[2pt]\scriptsize calcula la nota};
  \node[draw=guinda, line width=1.5pt, rounded corners=8pt, fill=white,
        text width=2.7cm, align=center, minimum height=1.5cm, right=1.6cm of server] (moodle)
        {Moodle\\[2pt]\scriptsize la nota aparece en el curso};
  \draw[-{Stealth}, line width=2pt, azulviv] (juego) -- (server)
        node[midway, above, font=\scriptsize\demi, text=azulviv]{aciertos};
  \draw[-{Stealth}, line width=2pt, guinda] (server) -- (moodle)
        node[midway, above, font=\scriptsize\demi, text=guinda]{AGS};
\end{tikzpicture}
\end{center}
\vspace{0.5cm}
\begin{center}
\begin{calloutbox}{verde}
\demi El maestro \textbf{no captura nada}. La calificaci\'on del libro de notas de
Moodle se actualiza autom\'aticamente al terminar la sesi\'on.
\end{calloutbox}
\end{center}
\end{frame}
\note{
\textbf{El cierre del c\'irculo (1 min).} Esta es la magia que m\'as impresiona.
Explica el flujo de derecha: el juego reporta los aciertos al servidor, el servidor
calcula la nota y la manda a Moodle usando AGS (el mecanismo de calificaciones de
LTI). Resultado: la nota aparece en el libro de calificaciones SIN que nadie la
escriba. Esta es la diapositiva para hacer una pausa y dejar que cale.
}

% ---------------------------------------------------------------------------
%  DIVISOR 02
% ---------------------------------------------------------------------------
\divisor{02}{C\'omo est\'a hecho}

% ---------------------------------------------------------------------------
%  ARQUITECTURA
% ---------------------------------------------------------------------------
\begin{frame}{La arquitectura, a vista de p\'ajaro}
\vspace{0.3cm}
\begin{center}
\begin{tikzpicture}[every node/.style={font=\demi\small}]
  \node[draw=guinda, line width=1.6pt, rounded corners=9pt, fill=guinda!6,
        text width=2.8cm, align=center, minimum height=1.4cm] (moodle)
        {\color{guinda}Moodle\\[1pt]\scriptsize plataforma del curso};
  \node[draw=azultec, line width=1.6pt, rounded corners=9pt, fill=azultec!6,
        text width=3.2cm, align=center, minimum height=1.4cm, right=2.0cm of moodle] (back)
        {\color{azultec}Servidor (Laravel)\\[1pt]\scriptsize la l\'ogica + la seguridad};
  \node[draw=azulviv, line width=1.6pt, rounded corners=9pt, fill=azulviv!6,
        text width=2.8cm, align=center, minimum height=1.4cm, right=2.0cm of back] (juego)
        {\color{azulviv}Juego (Unreal)\\[1pt]\scriptsize la pr\'actica};
  \node[draw=tinta!55, line width=1.4pt, rounded corners=9pt, fill=tinta!4,
        text width=3.2cm, align=center, minimum height=0.85cm, below=0.8cm of back] (db)
        {\color{tinta}Base de datos\\[1pt]\scriptsize alumnos, sesiones, notas};

  \draw[{Stealth}-{Stealth}, line width=1.6pt, guinda] (moodle) -- (back)
        node[midway, above, font=\scriptsize\demi, text=guinda]{LTI};
  \draw[{Stealth}-{Stealth}, line width=1.6pt, azulviv] (back) -- (juego)
        node[midway, above, font=\scriptsize\demi, text=azulviv]{enlace + API};
  \draw[-{Stealth}, line width=1.4pt, tinta!55] (back) -- (db);
\end{tikzpicture}
\end{center}
\vspace{0.25cm}
\begin{itemize}\small\setlength\itemsep{2pt}
  \item \hl{Moodle} manda al alumno y recibe la nota (v\'ia LTI).
  \item El \hlb{servidor} es el cerebro: valida, genera el enlace y guarda todo.
  \item El \textbf{juego} solo se enfoca en la pr\'actica y reporta resultados.
\end{itemize}
\end{frame}
\note{
\textbf{Arquitectura (1.5 min).} Tres piezas, f\'acil de recordar. (1) Moodle:
lo que ya tienen. (2) El servidor en Laravel: el cerebro; aqu\'i vive la seguridad,
la l\'ogica y la base de datos. (3) El juego en Unreal: solo se preocupa por la
pr\'actica. Las flechas: Moodle y servidor hablan por LTI; servidor y juego por un
enlace seguro + API. No entres en detalle t\'ecnico salvo que pregunten.
}

% ---------------------------------------------------------------------------
%  EL ENLACE SEGURO (MAGIC LINK)
% ---------------------------------------------------------------------------
\begin{frame}{El enlace que abre el juego es seguro}
\vspace{0.25cm}
\begin{columns}[T]
\column{0.5\textwidth}
Cuando el alumno abre la actividad, el servidor crea un \hl{enlace m\'agico} con un
token \'unico:
\begin{itemize}\setlength\itemsep{5pt}
  \item \demi De un solo uso: una vez abierto, \hl{ya no sirve}.
  \item \demi Con caducidad: vence solo (por defecto 2 horas).
  \item \demi Personal: identifica a ese alumno, ese curso, esa pr\'actica.
\end{itemize}
\column{0.5\textwidth}
\begin{tarjetabox}{azultec}
{\demi\color{azultec}?`Por qu\'e importa?}\\[3pt]
Nadie puede ``reusar'' el enlace de otro ni hacer trampa abriendo el juego por su
cuenta. La identidad viene \textbf{firmada} desde Moodle.
\end{tarjetabox}
\begin{calloutbox}{ambar}
\demi As\'i, la calificaci\'on siempre se le asigna a \textbf{la persona correcta}.
\end{calloutbox}
\end{columns}
\end{frame}
\note{
\textbf{Seguridad del lanzamiento (1 min).} El concepto a transmitir: el enlace que
abre el juego es como un boleto de un solo uso, con caducidad y a nombre del alumno.
Eso evita trampas (no puedes reusar el enlace de un compa\~nero) y garantiza que la
nota se le ponga a quien de verdad jug\'o. La identidad viene firmada desde Moodle,
no la inventa el juego.
}

% ---------------------------------------------------------------------------
%  PANEL DE MAESTROS
% ---------------------------------------------------------------------------
\begin{frame}{El panel para maestros}
\vspace{0.25cm}
\begin{columns}[T]
\column{0.5\textwidth}
\begin{tarjetabox}{guinda}
{\demi\color{guinda}Sin tocar Moodle ni c\'odigo}\\[3pt]
El maestro entra con un enlace propio y administra sus grupos y pr\'acticas desde
una pantalla sencilla.
\end{tarjetabox}
\begin{tarjetabox}{azulviv}
{\demi\color{azulviv}Compatible con LTI}\\[3pt]
El panel y la integraci\'on con Moodle \textbf{conviven}: lo que se hace en uno se
refleja en el otro.
\end{tarjetabox}
\column{0.5\textwidth}
Desde el panel puede:
\begin{itemize}\setlength\itemsep{4pt}
  \item Ver sus \hl{grupos} y a los alumnos.
  \item Generar \hl{enlaces} de pr\'actica.
  \item Consultar \hlb{resultados} y avance.
\end{itemize}
\end{columns}
\end{frame}
\note{
\textbf{Panel docente (1 min).} Aclara la duda t\'ipica: ``?`y si el maestro no usa
Moodle para todo?'' Tiene un panel propio, con enlace de acceso, donde ve grupos,
genera enlaces de pr\'actica y consulta resultados. Lo importante: el panel y LTI
NO se pelean, conviven. Es flexibilidad para distintos estilos de docente.
}

% ---------------------------------------------------------------------------
%  SEGURIDAD
% ---------------------------------------------------------------------------
\begin{frame}{Pensado con seguridad desde el inicio}
\vspace{0.3cm}
\begin{columns}[T]
\column{0.5\textwidth}
\begin{tarjetabox}{azultec}
{\demi\color{azultec}Identidad firmada}\\[2pt]\scriptsize
La identidad del alumno viaja firmada desde Moodle (LTI 1.3). No se puede falsificar.
\end{tarjetabox}
\begin{tarjetabox}{guinda}
{\demi\color{guinda}Enlaces de un solo uso}\\[2pt]\scriptsize
Tokens \'unicos, con caducidad, que se invalidan al usarse.
\end{tarjetabox}
\column{0.5\textwidth}
\begin{tarjetabox}{verde}
{\demi\color{verde}Conexi\'on cifrada (HTTPS)}\\[2pt]\scriptsize
Todo viaja cifrado, con certificado autom\'atico en el despliegue.
\end{tarjetabox}
\begin{tarjetabox}{azulviv}
{\demi\color{azulviv}Datos en casa}\\[2pt]\scriptsize
Cada escuela puede alojar su propio servidor y su propia base de datos.
\end{tarjetabox}
\end{columns}
\end{frame}
\note{
\textbf{Seguridad (45 s).} Cuatro mensajes, no m\'as: identidad firmada (no se
falsifica), enlaces de un solo uso, todo cifrado por HTTPS, y los datos se quedan en
la escuela. Si alguien de sistemas pregunta detalles, ofrece el documento t\'ecnico.
}

% ---------------------------------------------------------------------------
%  DIVISOR 03
% ---------------------------------------------------------------------------
\divisor{03}{Listo para usarse}

% ---------------------------------------------------------------------------
%  DESPLIEGUE
% ---------------------------------------------------------------------------
\begin{frame}{F\'acil de instalar y de replicar}
\vspace{0.25cm}
\begin{columns}[T]
\column{0.5\textwidth}
\begin{tarjetabox}{azulviv}
{\demi\color{azulviv}Un par de comandos}\\[3pt]
Se instala con Docker: la app, la base de datos y el HTTPS autom\'atico se levantan
juntos. No hay que configurar servidores a mano.
\end{tarjetabox}
\begin{tarjetabox}{verde}
{\demi\color{verde}En el servidor de la escuela}\\[3pt]
Pensado para vivir \textbf{on-premise}, junto a su propio Moodle, sin depender de la nube.
\end{tarjetabox}
\column{0.5\textwidth}
\begin{tarjetabox}{guinda}
{\demi\color{guinda}Un servidor por escuela}\\[3pt]
Si cada plantel quiere el suyo, se replica clonando el proyecto y cambiando un
archivo de configuraci\'on. Cada quien con sus datos.
\end{tarjetabox}
\begin{calloutbox}{ambar}
\demi Ya existe un \textbf{kit de despliegue} con gu\'ia paso a paso.
\end{calloutbox}
\end{columns}
\end{frame}
\note{
\textbf{Despliegue (1 min). } Mensaje para directivos/sistemas: instalar es f\'acil
(Docker, un par de comandos) y vive en el servidor de la escuela, junto a su Moodle,
sin nube obligatoria. Si quieren un servidor por plantel, se replica cambiando un
archivo. Ya hay un kit con gu\'ia. Esto baja el miedo de ``?`y qui\'en lo va a
mantener?''.
}

% ---------------------------------------------------------------------------
%  ESTADO ACTUAL
% ---------------------------------------------------------------------------
\begin{frame}{Qu\'e ya funciona hoy}
\vspace{0.35cm}
\newcommand{\ok}{\textcolor{verde}{\large\textbf{\checkmark}}\;}
\begin{columns}[T]
\column{0.5\textwidth}
\ok Lanzamiento desde Moodle real (LTI 1.3)\\[7pt]
\ok Selecci\'on de pr\'actica (Deep Linking)\\[7pt]
\ok Calificaci\'on de regreso a Moodle (AGS)\\[7pt]
\ok Enlaces seguros de un solo uso
\column{0.5\textwidth}
\ok Panel de maestros funcional\\[7pt]
\ok API documentada para el juego\\[7pt]
\ok Ejemplo de cliente en Unreal\\[7pt]
\ok Kit de despliegue + documentaci\'on
\end{columns}
\vspace{0.5cm}
\begin{center}
\begin{calloutbox}{verde}
\demi Todo el flujo se prob\'o de extremo a extremo contra un \textbf{Moodle real},
respaldado por \textbf{59 pruebas autom\'aticas}.
\end{calloutbox}
\end{center}
\end{frame}
\note{
\textbf{Prueba de que es real (1 min).} Esta es la diapositiva de credibilidad. No
es una maqueta: el flujo completo se prob\'o contra un Moodle real instalado para
eso. Menciona las 59 pruebas autom\'aticas como se\~nal de calidad. Recorre las
palomitas r\'apido; lo importante es la frase verde de abajo.
}

% ---------------------------------------------------------------------------
%  LA DEMO DE HOY
% ---------------------------------------------------------------------------
\begin{frame}{La demo de hoy}
\vspace{0.3cm}
{\medium\color{gristenue}\large Lo que van a ver, en orden:}
\vspace{0.45cm}

\begin{enumerate}\setlength\itemsep{8pt}
  \item[\bolanum{azulviv}{1}] Entrar a Moodle y abrir la actividad ``Pr\'actica en metaverso''.
  \item[\bolanum{guinda}{2}] El juego se abre solo con el enlace seguro.
  \item[\bolanum{azulviv}{3}] Realizar la pr\'actica (responder / acertar).
  \item[\bolanum{verde}{4}] Volver a Moodle y ver la \hl{calificaci\'on ya puesta}.
\end{enumerate}
\end{frame}
\note{
\textbf{Gu\'ia de la demo (clave, ya que el autor no est\'a).} Sigue EXACTAMENTE
estos 4 pasos en pantalla. (1) Inicia sesi\'on en Moodle con la cuenta de alumno de
prueba y abre el curso; ah\'i est\'a la actividad. (2) Dale clic: el juego se abre.
(3) Juega un poco para generar aciertos y termina la sesi\'on. (4) Regresa al libro
de calificaciones de Moodle y muestra que la nota ya aparece.
\medskip
\textbf{Si algo falla:} no improvises con la consola. Ten a la mano las capturas de
respaldo (carpeta de la demo) y di ``as\'i se ve cuando corre''. El soporte t\'ecnico
de Daniel queda para preguntas por correo.
}

% ---------------------------------------------------------------------------
%  CIERRE
% ---------------------------------------------------------------------------
{
\setbeamercolor{background canvas}{bg=azultec}
\begin{frame}[plain]
\begin{tikzpicture}[remember picture,overlay]
  \fill[guinda] ([yshift=2.6cm]current page.south west)
    -- (current page.south west) -- (current page.south east)
    -- ([yshift=1.1cm]current page.south east) -- cycle;
  \fill[ambar]  ([yshift=2.6cm]current page.south west)
    -- ([yshift=2.4cm]current page.south west)
    -- ([yshift=0.9cm]current page.south east)
    -- ([yshift=1.1cm]current page.south east) -- cycle;
  \node[anchor=west, text=white, font=\heavy\fontsize{34}{38}\selectfont, text width=12cm]
    at ([xshift=1.1cm,yshift=-2.6cm]current page.north west)
    {Gracias.};
  \node[anchor=west, text=white, font=\medium\fontsize{17}{22}\selectfont, text width=12cm]
    at ([xshift=1.15cm,yshift=-4.2cm]current page.north west)
    {El alumno juega.\\ La calificaci\'on llega sola.};
  \node[anchor=west, text=ambar, font=\demi\normalsize]
    at ([xshift=1.18cm,yshift=1.55cm]current page.south west)
    {Daniel Octavio Ram\'irez Neri \, $\cdot$ \, Metaverso Escolar TecNM $\times$ ITCJ};
\end{tikzpicture}
\end{frame}
}
\note{
\textbf{Cierre (30 s).} Repite el mantra: ``el alumno juega, la calificaci\'on llega
sola''. Invita a preguntas. Para dudas t\'ecnicas profundas, ofrece el documento de
dise\~no t\'ecnico y el correo de Daniel. Agradece el tiempo.
}
